\begin{abstract}

3D Gaussian Splatting (3DGS) represents a scene with explicit anisotropic Gaussian primitives and renders novel views through differentiable rasterization. Although this representation delivers high visual quality, standard 3DGS training remains expensive because it repeatedly grows and prunes a large Gaussian set, rasterizes many low-contribution Gaussian-tile pairs, and updates millions of Gaussian parameters through a general-purpose optimizer path.

This thesis studies how to accelerate single-scene 3DGS training while preserving the standard Gaussian representation, spherical-harmonic color model, front-to-back alpha blending semantics, and RGB photometric objective. The central observation is that training time can be viewed as the product of three factors: the average number of Gaussians during optimization, the average per-Gaussian cost of one training step, and the number of training iterations. Based on this decomposition, the proposed system, RapidGS, reduces redundant Gaussian growth through multi-view structure control, lowers the training-step cost through a fused CUDA backend and compact rasterization, and uses a shorter training schedule under an explicit quality constraint.

The current main pipeline uses COLMAP/SfM initialization and an RGB-only training objective. Optional Gaussian initialization from AnySplat is treated as a separate initialization path rather than a required component of the main result. On the Mip-NeRF 360 nine-scene benchmark under the current local evaluation protocol, RapidGS reduces average training time from 1053.54 seconds for the original 3DGS baseline to 152.47 seconds, while improving PSNR, SSIM, and LPIPS. Compared with a local FastGS-quality reference, it reduces training time by approximately 35.4 percent while maintaining or improving reconstruction quality.

\end{abstract}
